\documentclass[11pt]{article}

\usepackage[a4paper,margin=1in]{geometry}
\usepackage{booktabs}
\usepackage{xurl}
\usepackage{array}
\usepackage{microtype}
\usepackage{fancyhdr}
\usepackage[hidelinks]{hyperref}

\hypersetup{
  pdftitle={Model Intervention Records: A Conformance Test for Weight-Space Claims},
  pdfsubject={MIR-1: conformance requirements and test suite},
  pdfkeywords={weight-space programming, model intervention records, weight
    editing, conformance testing, reproducibility, evaluation validity}
}

\setlength{\parskip}{0.35em}
\setlength{\parindent}{0pt}

\pagestyle{fancy}
\fancyhf{}
\fancyhead[L]{\small\textsc{mir-1}}
\fancyhead[R]{\small Model Intervention Records}
\fancyfoot[C]{\small\thepage}
\renewcommand{\headrulewidth}{0.4pt}
\renewcommand{\footrulewidth}{0pt}

\begin{document}

\thispagestyle{plain}
\begin{center}
\vspace*{0.4em}
{\LARGE\bfseries Model Intervention Records}\\[0.35em]
{\LARGE\bfseries A Conformance Test for Weight-Space Claims}\\[0.8em]
{\large Version 1.0}\\[0.25em]
{\large 14 November 2026}\\[0.8em]
{\small\textbf{Keywords:} weight-space programming; model intervention records;
weight editing; conformance testing; reproducibility; evaluation validity}
\end{center}
\vspace{0.2em}
\hrule
\vspace{0.6em}

\begin{abstract}
Claims that an intervention changes a language model's behavior are currently
not comparable. The interventions at issue are the ones now described as
weight-space programming: a model's weights read and edited directly, without
gradient descent or a training loop. Two reports of the same intervention can
both be accurate and still disagree about what happened, because neither record
states what was required not to change. We specify a record format with three
mandatory components: an executable falsifier; an invariant set containing at
least one metric outside the claim's own declared scoring function; and
content-addressed lineage to a base checkpoint. We then specify a mechanical
test that checks them, and we apply it to a published intervention: the record
satisfies ten of eleven checks and fails the one on cost disclosure. The test
emits no aggregate score and does not assess importance, safety, or desirability.
We release two dependency-free implementations, written from the same text in
different languages, and they agree on every input in the repository. Adoption
requires no authority: conformance is a claim that anyone may check.
\end{abstract}

\section{Introduction}

Weight-space programming---reading and writing the weights of a trained model
directly, without gradient descent or a training loop---is now cheap enough to
be routine. A search loop can propose an intervention, score it on a small set,
keep it only if the gain exceeds measured noise, and repeat, at a cost within
reach of a single researcher with consumer hardware. Reports of such
interventions appear within days of the methods that produce them, and the
artifacts---modified checkpoints, run logs, rejected candidates---are published
alongside.

The reporting format has not kept pace. A report typically states a base model,
an intervention, a target metric, and a result. That set of fields is
sufficient to describe a success. It is not sufficient to describe a success
that damaged something else, because there is no field in which the author can
state what was required not to change.

The consequence is not that authors are careless. It is that the format
determines the observation set. When the scoring function is the only
specification of what to measure, the record will describe exactly what was
scored. Two authors using the same format will produce records that are
individually accurate and jointly incomparable, and neither record will contain
the information needed to detect that the intervention broke a property nobody
was looking at.

We make five contributions:

\begin{enumerate}
\item A documented instance of this failure, from a public run on 8 November
  2026~\cite{beaker}, in which an accepted intervention reduced mathematical
  accuracy from 80\% to between 23\% and 54\% and GSM8K accuracy from 46\% to
  0\% in a pre-check, while the score the loop optimized stayed accurate.
\item A record format that requires the missing field, together with the fields
  needed to make a claim checkable by a third party. We call it MIR-1, for
  Model Intervention Records.
\item A mechanical conformance test over that format, with an explicit list of
  what it does not check.
\item Two reference implementations, dependency-free, running on one machine.
\item A worked example: a published intervention re-recorded under the format,
  with the three fields the exercise showed were missing and one defect it
  exposed in the format.
\end{enumerate}

We do not claim that weight editing is the correct unit of behavioral change,
and we do not claim that adopting this format makes a model safer. Both
questions are out of scope, and Section~\ref{sec:cannot} explains why leaving
them out is a condition of the format being usable at all.

\section{The failure mode}

On 8 November 2026 a public run reported an intervention that reduced response
length~\cite{beaker}. The selection score was brevity credit, awarded only to
answers that were right: an edit was kept when it shortened answers without
making them wrong on the target set. The loop accepted one, and the score
described the reduction accurately.

The run's write-up reports what the score did not. Mathematical accuracy fell
from 80\% to between 23\% and 54\%; GSM8K accuracy fell from 46\% to 0\% in a
pre-check; at least one run began every mathematical response with a stray
period. The author noticed, and said so in prose, after the loop had finished.

Nothing in that sequence is anomalous. The loop did what loops do: it optimized
the function it was given, over the observations the format could express. The
score was accurate and the system was damaged, and both statements hold at once
because the score never claimed otherwise---it had no field in which to.

The prose is the part that does not travel. It is not read by the next loop, it
is checked by nothing, and a reader who wants the number has to find the
sentence that contains it.

Three properties of this failure make it a format problem rather than an author
problem.

First, it is reproducible by anyone. The same field set, applied to a different
intervention, produces the same class of blindness. The failure does not depend
on the author's judgment, care, or honesty.

Second, it scales with the loop. A manual intervention is inspected once by a
human who may notice a side effect. A loop runs thousands of interventions and
keeps the ones that pass a score. The score is the entire specification of what
to keep, so the blindness is applied thousands of times at machine speed.

Third, it is invisible in the output. A damaged model that satisfies its score
looks like a success in every field the record contains. Detection requires an
observation the record did not request.

The remedy follows directly: the record must contain a field for what was
required not to change, and that field must be checked mechanically, because a
field that is checked by inspection is a field that is checked when someone
remembers.

\section{Record requirements}

The requirements below use MUST in the sense of RFC~2119~\cite{rfc2119}. A
record that omits a MUST is not a MIR-1 record, whatever else it contains.

\subsection{Claims as tuples}

A claim is a tuple of behavior, distribution, direction, threshold, and the
metrics the scoring function observes. The
behavior must be defined operationally by the harness. A behavior defined in
prose is not checkable, and a claim that is not checkable cannot be the subject
of a conformance test.

\subsection{Falsifier}

A record must carry an executable falsifier: a command or a fully specified
measurement that returns a boolean. The requirement is not that the author be
wrong; it is that the author state, in a form a machine can evaluate, what
observation would show that they are.

A falsifier written in prose does not conform, regardless of its precision. The
distinction is not stylistic. Prose falsifiers are read by people who already
believe the claim; executable ones are run by people who do not.

\subsection{Invariant set}

An invariant set is a list of properties, each naming a metric, a direction, and
a numeric bound. The direction says how the bound is read: a floor, a ceiling,
or a signed change from the base run. The claim, in turn, declares the metrics
its scoring function observes in \texttt{scoring\_metrics}, and its own behavior
must be one of them. The set must contain at least one metric that is not among
them.

This is the requirement that addresses Section~2. A record whose invariant set
is a subset of its scoring function reproduces the failure exactly while
presenting every required field as complete---which is worse than omitting the
field, because it produces the appearance of diligence.

The minimum form is one metric, measured outside the scoring function, with a
bound. In the case of Section~2 the set would have contained mathematical
accuracy with a bound of five percentage points: the number the prose reported
would have been in a field a machine can read, next to the claim it qualifies,
rather than in a paragraph below it.

The check is on the declaration, but the comparison is mechanical. The suite
does not reconstruct the scoring function from the harness; it compares each
invariant's metric against the metrics the claim declares. The degenerate case
is therefore closed: an invariant whose metric is the claim's own behavior is
inside by construction, and no field can place it outside. What remains open is
the general false declaration---a record that scores a metric and omits it from
its declared set will pass C-05, and that is a false record, not a
non-conforming one. The suite can make the honest omission visible, which is the
failure in Section~2. It cannot make a false declaration true, and it does not
claim to.

\subsection{Lineage and content addressing}

Every artifact referenced by a record is identified by a hash of its bytes: the
base checkpoint, the intervention payload, the harness, and each run's inputs.
Lineage is a directed edge from a record to the record or records it was
applied on top of, resolving to a base artifact.

Content addressing does two things that identifiers do not. It makes silent
modification detectable, and it makes a fork legible: two records that share a
parent are visibly related, and a record whose parent cannot be retrieved is
visibly orphaned.

A hash that resolves to content nobody retained is the common failure of
reproducibility, and it is a failure of the record rather than of the reader.
The requirement is retrieval by a third party, not the existence of a digest.
Because a checkpoint is frequently larger than anything a reviewer will
download to check a record, the format allows a locator: a file named for the
digest that declares the size and a URL, and is confirmed against the host by a
request that transfers no bytes. The check reports whether it confirmed the
locator or merely accepted it.

\subsection{Population}

Every run declares how many units it measured and how many samples it drew per
unit. This is the field whose absence makes two accurate records incomparable,
and Section~\ref{sec:worked} shows a published case in which it is the
difference between a 46-point improvement and two different experiments.

\subsection{Independence}

A claim is replicated when at least one run was executed by an operator other
than the claimant, whose declared funding source is disclosed and differs from
the claimant's. The record names the claimant: a replication is measured
against whoever asserts the claim. A funding source that is not disclosed is not
a source, so a record declaring one cannot reach the status. The mechanism
for establishing distinctness is not specified here, and this is a known gap:
funding disclosure is a minimum, not a solution.

\subsection{Export}

A registry must be exportable in full, by one command, in a format a third
party can re-import and verify.

This requirement protects the people whose records the registry holds, and it
is the only requirement in the format that does so. A registry that cannot be
left is not a shared record. It is a silo with a public address, and the
distinction is the difference between a format and a product.

\section{The test}

The suite implements eleven checks: parse and content-address verification;
the published schema, run rather than restated; hash resolution as stored
bytes or as a confirmed locator; falsifier declaration and return type;
invariant metrics, directions, and numeric bounds, with at least one metric
outside the claim's declared scoring function; lineage
resolution; cost and wall time disclosed as numbers on every run, where a run
that produced a result must have taken time; manifest
resolution; an export round-trip that must reproduce its own results; a
status check, which refuses \texttt{replicated} unless the record names a claimant with
a disclosed funding source and carries the number of runs the manifest requires
by other operators with disclosed, differing sources; and a grounding check,
which requires every run to report the metric the claim is about.

Two design constraints govern the check set.

The first is implementability. The suite must be implementable from the
specification alone, within one working day, by an implementer with no access
to the authors. A requirement that fails this condition does not enter version
1. This is a quality gate on the specification rather than on the
implementation, and it is the reason the requirement list is short.

The second is the exclusion of non-mechanical checks. The suite does not report
on importance, novelty, evaluation quality, operator reputation, model safety,
or desirability. These are not mechanical, and a suite that reports on them
becomes a rating. A rating has an owner, an owner has an interest, and an
interest is the thing that the format exists to prevent from entering the
record.

For the same reason the suite emits no aggregate score. An aggregate is
optimized rather than satisfied, and a score that can be optimized by damaging
an unmeasured property reproduces Section~2 at a higher level of abstraction.

The falsifier is checked for declaration and return type rather than executed
by default. A record is untrusted input, and a validator that executes commands
from untrusted input cannot be run on the records it is meant to check.
Execution is available as an explicit option for records whose provenance is
known.

Each record-level check except C-07 and C-09 has a mutation that makes it fail;
C-07's failure is the record as published, and C-09 is the serializer's
round-trip, which no record mutation can remove. A check that has never been
observed to fail is not known to work, and the repository carries the cases that
produce each failure it can produce.

The schema check is worth one sentence about how it is built. The validator
does not carry its own list of required fields: it loads the published schema
and checks the record against it. It did carry one, and the copy drifted --- a
field became required in the schema and stayed missing from both
implementations, which agreed with each other and disagreed with the document.
Two implementations agreeing proves the text is implementable; it does not
prove either one matches the contract.

The content address in the first check forced a normative detail that a prose
specification usually leaves out: the canonical bytes of a record. Writing the
second implementation from the text is what made it normative, because two
serializers that spell the same value differently produce two ids for one
record, and a format whose identity depends on a serializer cannot be
implemented twice. The detail that had to be fixed was number formatting: the
first implementation wrote exponent notation for values at or above $10^{16}$
and below $10^{-4}$, while canonical bytes may not contain an exponent. The
text said the canonical form is plain decimal; the implementation did not match
the text. It does now. This is the class of defect a second implementation is
for, and it was found before the format was used by anyone else.

\section{The worked example}
\label{sec:worked}

We re-recorded a published intervention under the format. On 2 November 2026 an
LLM changed 131,072 weights of Qwen3.5-0.8B, 0.017\% of the model, to remove a
reasoning loop that made it repeat until the token limit. Two operators
published measurements: the original group, and a rebuild from the published
edit by a second operator that re-ran the evaluation.

\begin{table}[h]
\centering
\begin{tabular}{@{}l p{0.76\textwidth}@{}}
\toprule
Check & Result on the record \\
\midrule
C-01, C-02 & pass \\
C-03 & pass: 21 hashes resolve; 8 references by locator, confirmed against the host with \texttt{--verify-remote} \\
C-04 & pass: declared; executed with \texttt{--execute-falsifier}, returned true, recomputes 0.51 from 2,000 samples \\
C-05 & pass: 3 invariants outside the declared scoring function; bounds declared, not evaluated \\
C-06 & pass: applies directly to a base artifact \\
C-07 & fail: no run states a cost and wall time \\
C-08, C-09 & pass \\
C-10 & pass: status \texttt{submitted}; lineage resolves \\
C-11 & pass: all 4 runs report \texttt{math500\_accuracy} \\
\bottomrule
\end{tabular}
\caption{The eleven checks applied to a published intervention re-recorded under
MIR-1. The record is `submitted', not `replicated', because of the single
failure.}
\end{table}

The claim is that the edited checkpoint reaches at least 45\% on MATH-500. The
original operators report 48.1\% over 500 problems with 4 samples each; the
rebuild by a second operator reports 51.0\% over the same population. Two
operators, one re-run that lands inside the claim. The record's status is \texttt{submitted}
regardless, for two separate reasons: no run states a cost, so no run counts
toward any requirement; and the claimant's funding source is not disclosed, so
the status check cannot establish an independent source either. That is the
format declining to credit evidence it cannot see priced. The requirement costs
one number per run.

Three fields were added because the record could not be written without them.

A baseline run applies no intervention, so its intervention hash is null. A
change that is not measured against a baseline cannot be stated as a change.

Cost and wall time may be null, meaning explicitly undisclosed. The field stays
in the record and the check still fails, so the omission is visible in the
record rather than absent from it.

A checkpoint of 1.7\,GB will not be downloaded to check a record, so a hash may
resolve to a locator, and the check states whether it confirmed the locator
against the host. A validator that reports an unconfirmed locator as resolved
is claiming more than it measured.

The example also exposed a defect we did not fix. The model repository contains
two files named for the same distribution. One measures 25 problems and reports
5.0\%; the other measures 500 and reports 51.0\%. The model card states that the
first is a smoke test and the full baseline is in a different repository, under
a different commit; the file name does not state that, and neither file carries
a population. A reader who computes the difference between the two files has
computed a 46-point improvement across two populations. With
\texttt{n\_units} required, the difference sits inside the record where it can
be seen. Nothing prevents the comparison, and version 1 checks only that every
run states its population, not that two runs share one.

\section{What the test cannot do}
\label{sec:cannot}

Three threats are mitigated by the format: evaluation capture, because the
invariant set requires an observation outside the scoring function; record
spam, because a record costs a run; and a comparison between two populations,
because every run states its own.

Three are not. Sybil replication is only partially addressed by funding
disclosure. Harness decay is reduced but not eliminated by the retrieval
requirement. And a well-funded actor holding the only credible evaluations can
satisfy every requirement in the format while controlling what counts as
evidence.

One limit sits inside the record. The invariant field is verified against the
claim's own declaration of its scoring function: a record that scores a metric
and omits it from that declaration passes C-05, because the suite cannot see the
scoring function itself. The limit produces a false record rather than an
unparseable one. The suite makes the honest omission visible, which is the
failure in Section~2; it does not detect a declaration that is untrue.

The status vocabulary is restricted rather than checked loosely. Four statuses
exist, and each is a function of evidence: a record whose lineage does not
resolve must be \texttt{orphan}, and one whose lineage resolves must not be; \texttt{replicated}
requires a disclosed claimant and the manifest's independent runs. The three
statuses a registry would normally hold --- \texttt{challenged}, \texttt{refuted}, \texttt{expired} ---
are not in the format, because each needs an object the format does not define.
A status no check can test is a word, and the vocabulary is limited to what the
evidence determines.

We state the last one explicitly because a threat model that lists only
mitigated threats is a prospectus. The format does not fix this condition. It
requires that evaluations be named, hashed, and retrievable, which makes the
control observable without making it removable. Observability is a smaller
claim than safety, and it is the only claim the test can support.

The worked example also shows what the test does not reach. The dataset used by
the harness is loaded without a revision pin, so the harness hash covers the
code and not the data it ran on. A record can therefore be conformant and still
be tied to a distribution that has since changed. Pinning the dataset is a
change to the harness, not to the suite, and version 1 does not require it.

Two further gaps were found by asking what a published artifact of a different
kind would look like as a record, and both are reported here rather than
silently excluded from the check list.

The first is that lineage is declared, not computed. Every hash must resolve,
and no check asks whether the payload was derived from the base, or whether it
differs from it. The repository now carries the case: set the payload hash equal
to the base hash, keep the claim, and the record fails exactly the check the
honest record fails and nothing else. A record can assert that it edited a
checkpoint into something else while the something else is the checkpoint. The
test cannot see this, because the test asks whether a hash resolves and not
whether a transformation happened. Closing it requires the base bytes and a
definition of what counts as derived, which is a larger requirement than
resolving a hash; version 1 states the gap instead of pretending to close it.

The second is that a construction cannot be typed. The format requires a base
artifact that resolves, so an intervention that builds a model without a parent
--- weights written directly, no training, no base to edit --- has no value for
that field. Such work is out of scope in version 1: not rejected, not expressed.
This is a scope limit rather than a defect in the checks, and it matters because
the artifacts the field is producing are increasingly of that kind.

\section{Related work}

Reporting checklists have solved an adjacent problem in medicine.
CONSORT~\cite{consort} specifies the minimum information a randomized trial
must report, and became effectively mandatory through editorial policy rather
than through any enforcement authority. MIQE~\cite{miqe} did the same for
quantitative PCR, where the claim being reported is a measurement rather than
an outcome. In both cases the mechanism is the same: a venue requires the
checklist, and the checklist is short enough that requiring it costs authors
little.

Supply-chain attestation addresses a different question with a similar
structure. SLSA~\cite{slsa} specifies what must be provable about how an
artifact was built, and it is enforced by tooling rather than by inspection.

Model Cards~\cite{modelcards} and Datasheets~\cite{datasheets} describe
artifacts: a model, a dataset. The object of this format is different. It
describes an intervention, which has a parent, and the claim is about a change
rather than a thing. A model card states what a model is. This states what
someone did to it, what followed, and what was required not to follow.

Two recent lines of work supply the objects this format is meant to describe.
Weight-space programming~\cite{sparks} reports that a general-purpose model can
read and write the weights of another network directly, and
WeightLoop~\cite{weightloop} turns that into a loop driven by experimental
feedback. Deletion Tuning~\cite{deletion} reaches comparable results by masking
a small fraction of weights during post-training, with a much smaller artifact
and less forgetting. Each is a report of an intervention, and each would be
comparable to the others under a shared record.

The authors of weight-space programming have since stated that it is not a
method but ``the name of an emerging concept'', and that their experiments are
``evidence for that statement, not a proposed method for editing or training
networks''~\cite{sparks}. They also report more than a hundred replications,
extensions, benchmarks, and applications known to them. This is the condition
the format addresses: a concept with many independent reports and no shared
record yields many claims that cannot be compared, and the missing piece is not a
method specification but the record.

The \texttt{dspy-wsp} package turns the concept into a DSPy
optimizer~\cite{dspy}: a frontier model reads and edits a student's weights with
no gradients, and the result is a program that can be optimized like any other.
Runs that accept an edit carry a \texttt{patch/manifest.json} naming the base
model and revision, the writer, the forward-pass budget, the changed tensors, and
the result of the final check on the 400-case test split, with the tensor delta
in \texttt{edit.npz} beside it. The run report carries the budget, the turn
count, every validation run, and a patch summary. Where cost and wall time are
published, they appear in the results README rather than in the report, which is
the run-level gap C-07 measures. The distance is small: the format is not asking
for a new logging system, only for two numbers the operator already has. The
edit in the d1 run is a 1,024-weight correction to one scorer layer, computed in
closed form; the point here is the artifact it leaves, not the search that
produced it.

Interpretable weight-space programming~\cite{iwsp} is the closest work and the
sharpest contrast. An agent writes a transformer's weights by hand, with no
gradient descent, seeing one validation score per iteration, and the resulting
model is small enough to be read; a fresh model rebuilds it from a two-hundred
word description, and thirteen of thirteen rebuilt models reach full accuracy on
held-out algorithmic tasks. That is a stronger property than anything this
format checks. A record asks whether an artifact can be retrieved; iWSP asks
whether it can be regenerated from a description, and it can.

The two answers are not rivals, and the scope of each follows from the artifact.
Where the artifact fits in a page, interpretability is the better control,
because a reader can see what the model computes and no record format can
substitute for that. Where the artifact is an edit to a checkpoint of a size
nobody will read, readability is not available, and the only mechanical
controls left are the ones specified here: an invariant outside the scoring
function, a falsifier, and lineage that resolves. The iWSP runs also publish
\texttt{cost\_usd} and wall-clock seconds per session in machine-readable
metadata, which is precisely the number that C-07 requires and the one the
worked example lacks. The requirement is not exotic; it is already being met by
a project that is not writing records.

\section{Limitations}

No record of a published intervention in this paper conforms. The worked
example fails the check on cost disclosure; the example in Section~2 fails the
check on invariants. A synthetic record in the repository passes all eleven
checks, so the test is satisfiable; it has not been satisfied by a published
intervention without a change in what operators publish. That change is one
number per run. It has not been made in a record, and the reason is not that the
number is hard to produce: at least one published project writes it into
machine-readable metadata for every run it performs~\cite{iwsp}. The gap is that
the number is not written where a checker can read it, which is the only thing
this format asks.

The unit question is unresolved. It is not known whether a weight edit can
express a behavior that a prompt or a steering vector cannot, and the record
format is therefore typed: an intervention may be declared as weights, vector,
harness, or other. A record declaring the cheaper unit is as valid as one
declaring the more expensive one, and is often the truthful one. If it turns
out that the cheaper unit is always sufficient, the format remains correct and
the interesting objects become vectors and prompts.

The independence mechanism is underspecified, as noted in Section~3.6. Operator
identity is taken from what the record declares, and a record that declares a
false funding source is not detectable by the suite.

Invariant bounds are declared, not evaluated. The suite checks that a bound
exists, that it is numeric, that its direction is one of three values, and that
at least one invariant's metric lies outside the declared scoring function. It
does not compute whether the bound held. Satisfaction is a claim the record
makes, and a reader who wants it must run the harness the record names.

Two of the example records fail, on different checks: one on cost, and one
synthetic fixture on invariants. A format whose best available examples are failures
has been shown to be checkable; it has not been shown to be satisfied by a real
artifact.

No claim is made about adoption. Whether venues will require the checklist, and
whether registries will satisfy the export requirement, are empirical questions
that this paper cannot answer. The format is published with a stated failure
condition for that reason.

\section{Availability}

The specification, schema, validator, normative kernel, the worked example, and
the example records are released under CC BY 4.0 for text and Apache-2.0 for
code. Both validators have no third-party dependencies and run on a single
machine. The kernel is a standalone file of under 250 words, intended for
inclusion in automated review pipelines and in the context window of
model-based review systems; it is normative, and where it differs from the
specification the specification governs.

The worked example is reproducible without downloading either checkpoint: the
record cites locators, the validator confirms them against the host in one
request each, and the falsifier recomputes the claim from the per-sample
observations one of the operators published. From a copy of the repository:

\begin{verbatim}
python3 validate.py records/MIR-R1.json --verify-remote
python3 records/falsify.py records/MIR-R1.json
python3 crosscheck.py --remote
\end{verbatim}

The third command runs the reference implementation and the second
implementation on the same records, every mutation in the repository, a corpus
of damaged stores, and, with \texttt{--fuzz}, a set of random documents. It
exits non-zero if their reports differ by one byte.

The random corpus is half arbitrary JSON shapes and half mutations of a
conforming record, and the run prints, per check, how many documents passed it,
how many failed it, and how many distinct messages it produced. The distinction
matters: an arbitrary shape fails the schema and everything downstream
short-circuits, so a corpus of shapes alone yields one message per check and
measures agreement about malformed input. Over 400 documents the mixed corpus
reaches 190 distinct messages, and every check but the round-trip invariant has
been seen to fail.

A single implementation is a project. A second implementation is evidence that
the text, rather than its author, specifies the format. We have two:
\texttt{validate.py} in Python and \texttt{impl2/} in Go, written from the
specification and sharing no code. They agree on every input in the repository.
Writing the second one found the canonical-form defect described in
Section~4. Four further defects found while extending the differential test are
logged in the repository's defect notes, each with the test that now guards it.
That is evidence of implementability; it is not evidence of
implementer independence, because the same author wrote both. The condition for
a standard remains an implementation by someone with no access to the authors,
and the specification is written so that this is possible within one working
day.

\begin{thebibliography}{10}

\bibitem{consort} Schulz KF, Altman DG, Moher D. CONSORT 2010 statement:
  updated guidelines for reporting parallel group randomised trials. BMJ. 2010.

\bibitem{miqe} Bustin SA, et al. The MIQE guidelines: minimum information for
  publication of quantitative real-time PCR experiments. Clinical Chemistry.
  2009.

\bibitem{rfc2119} Bradner S. Key words for use in RFCs to indicate requirement
  levels. RFC 2119. IETF, 1997.

\bibitem{modelcards} Mitchell M, et al. Model cards for model reporting. FAT*
  2019.

\bibitem{datasheets} Gebru T, et al. Datasheets for datasets. Communications of
  the ACM, 2021.

\bibitem{slsa} OpenSSF. SLSA: Supply-chain Levels for Software Artifacts.
  Specification.

\bibitem{sparks} Takehi R, Teoh J, Zhang A, Khattab O. Sparks of Weight-Space
  Programming: LLMs can natively read \& write neural network weights. 2
  November 2026. Response to Comments added 12 November 2026.
  \url{https://rikiyat.github.io/blog/weight-space-programming/}

\bibitem{dspy} Breunig D. dspy-wsp: weight-space programming as a DSPy
  optimizer. Version 0.2.0, 13 November 2026.
  \url{https://github.com/cmpnd-ai/dspy-wsp}. MIT License.

\bibitem{hfmodel} RikiyaT. Qwen3.5-0.8B-weight-space-programming. Hugging Face
  model repository, 2 November 2026.

\bibitem{weightloop} Li M, Li B, Zhao W, Zhang C, Zhang B, Xu J, Zang Z, Leung
  VCM, Li SZ. WeightLoop: Weight-Space Programming with Experimental Feedback.
  arXiv:2611.07319, 2026.

\bibitem{beaker} Nayebi R. Adventures in Weight-Space Programming with Beaker.
  8 November 2026.
  \url{https://www.withbeaker.ai/blog/weight-space-programming}. Public run
  \texttt{3230a4ea}:
  \url{https://app.withbeaker.ai/public/runs/3230a4ea-512b-4a9f-b42a-4d271cca03fd}.

\bibitem{iwsp} Mehta Y, Singh C, Xu W, Antonello RJ, Flinker A, Mesgarani N,
  Bonner MF, Jojic N, Gao J. Interpretable weight-space programming builds
  transparent transformers and brain encoding models. 2026.
  \url{https://www.alphaxiv.org/abs/2611.interpretable-weight-space-programming}

\bibitem{deletion} Lu J, Yang Z, Lasby M, Pote T, Ioannou Y, Ren M. The
  Surprising Effectiveness of Deleting Weights in LLM Post-Training.
  arXiv:2611.10866, NeurIPS 2026.

\end{thebibliography}

\end{document}
